\section{Контекст}
\label{sec:context}

\texttt{PySATL Core} представляет собой вычислительное ядро экосистемы \texttt{PySATL}. Его назначение --- предоставить независимое от конкретных реализаций распределений API для:
\begin{itemize}[compact]
	\item создания распределений из параметрических семейств;
	\item вычисления функциональных и числовых характеристик по единым контрактам;
	\item генерации выборок с настраиваемой стратегией сэмплирования;
	\item расширения системы пользователем без модификации исходного кода ядра.
\end{itemize}

\textbf{Внешние акторы и окружение.}
Ядро используется конечными пользователями, которым требуется вычислять характеристики и строить выборки, и другими пакетами \texttt{PySATL}, которые строят на базе ядра прикладные методы математической статистики. На уровне инфраструктуры ядро опирается на стандартную среду Python и числовой стек (в первую очередь массивы \texttt{NumPy}); при этом архитектурной целью является уменьшение зависимости прикладных пакетов от \texttt{scipy.stats} за счёт переноса вычислительной модели в собственные модули.

\textbf{Граница системы.}
В рамках ядра фиксируются контракты распределения, семейств, поддержек (support), стратегий вычислений и сэмплирования, а также глобальный реестр преобразований характеристик. Вне границы системы остаются:
\begin{itemize}[compact]
	\item визуализация, статистические тесты и прикладные алгоритмы верхнего уровня;
	\item конкретные численные методы интегрирования/оптимизации, если они предоставляются внешними библиотеками;
	\item хранение данных и пользовательские интерфейсы прикладных пакетов.
\end{itemize}

\textbf{Точки расширения.}
Архитектура предполагает, что пользователи и пакеты экосистемы могут:
\begin{itemize}[compact]
	\item регистрировать новые семейства и параметризации;
	\item добавлять новые аналитические реализации характеристик;
	\item добавлять новые вычислительные методы и правила вывода в глобальный реестр характеристик;
	\item подменять стратегии вычислений и сэмплирования для конкретного распределения.
\end{itemize}
